\documentclass[../tesis.tex]{subfiles}
\begin{document}

\chapter{Clasificación de métodos de predicción}
\label{apendice:clasificacion_ml}

La siguiente tabla ilustra los métodos de aprendizaje automático revisados en la literatura,
clasificados según el \textbf{tipo de datos de entrada} ---Datos Clínicos Tabulares frente a
Imágenes/Radiómica (CT/DXA)--- para la \textbf{clasificación transversal del riesgo de osteoporosis}.
Esta clasificación justifica la escalera de algoritmos propuesta en el Capítulo~3 y evidencia
que el enfoque multimodal de \textit{DeepDXA} atiende la brecha de integración multimodal
menos explorada en el estado del arte latinoamericano.

\vspace{0.5em}

\begin{table}[h]
\centering
\caption{Clasificación de la literatura principal de clasificación transversal de osteoporosis según el tipo de datos de entrada.}
\label{tab:clasificacion_ml}
\renewcommand{\arraystretch}{1.6}
\rowcolors{2}{gray!12}{white}
\begin{tabular}{p{4.2cm} p{5.2cm} p{5.2cm}}
\toprule
\rowcolor{tableheader}
\textbf{Objetivo clínico} &
\textbf{Datos clínicos / Tabulares} &
\textbf{Imágenes / Radiómica (CT/DXA)} \\
\midrule
\textbf{Clasificación transversal} \newline
(Normal vs.\ Osteoporosis) &
\textit{Random Forest / XGBoost / SVM} \newline
Carvalho \& Gavaia~\cite{Carvalho2025} \newline
Je et al.~\cite{Je2025} \newline
Jin et al.~\cite{Jin2025} \newline
Yoo et al.~\cite{Yoo2013} \newline
Miranda et al.~\cite{Miranda2021} \newline
\textit{(SVM / BDAT ultrasonido, Chile)} &
\textit{CNN / Deep Learning} \newline
Ong et al.~\cite{Ong2023} \newline
Rahim et al.~\cite{Rahim2023} \\
\bottomrule
\end{tabular}
\end{table}

\clearpage
\end{document}
